\DocumentMetadata{
  pdfstandard={ua-2},
  tagging=on,
  lang=en,
  xmp=true}
% Options for packages loaded elsewhere
% Options for packages loaded elsewhere
\PassOptionsToPackage{unicode}{hyperref}
\PassOptionsToPackage{hyphens}{url}
\PassOptionsToPackage{dvipsnames,svgnames,x11names}{xcolor}
%
\documentclass[
  english,
  11pt,
]{article}
\usepackage{xcolor}
\usepackage[margin=1in]{geometry}
\usepackage{amsmath,amssymb}
\setcounter{secnumdepth}{-\maxdimen} % remove section numbering
\usepackage{iftex}
\ifPDFTeX
  \usepackage[T1]{fontenc}
  \usepackage[utf8]{inputenc}
  \usepackage{textcomp} % provide euro and other symbols
\else % if luatex or xetex
  \usepackage{unicode-math} % this also loads fontspec
  \defaultfontfeatures{Scale=MatchLowercase}
  \defaultfontfeatures[\rmfamily]{Ligatures=TeX,Scale=1}
\fi
\usepackage{lmodern}
\ifPDFTeX\else
  % xetex/luatex font selection
\fi
% Use upquote if available, for straight quotes in verbatim environments
\IfFileExists{upquote.sty}{\usepackage{upquote}}{}
\IfFileExists{microtype.sty}{% use microtype if available
  \usepackage[]{microtype}
  \UseMicrotypeSet[protrusion]{basicmath} % disable protrusion for tt fonts
}{}
\makeatletter
\@ifundefined{KOMAClassName}{% if non-KOMA class
  \IfFileExists{parskip.sty}{%
    \usepackage{parskip}
  }{% else
    \setlength{\parindent}{0pt}
    \setlength{\parskip}{6pt plus 2pt minus 1pt}}
}{% if KOMA class
  \KOMAoptions{parskip=half}}
\makeatother
% Make \paragraph and \subparagraph free-standing
\makeatletter
\ifx\paragraph\undefined\else
  \let\oldparagraph\paragraph
  \renewcommand{\paragraph}{
    \@ifstar
      \xxxParagraphStar
      \xxxParagraphNoStar
  }
  \newcommand{\xxxParagraphStar}[1]{\oldparagraph*{#1}\mbox{}}
  \newcommand{\xxxParagraphNoStar}[1]{\oldparagraph{#1}\mbox{}}
\fi
\ifx\subparagraph\undefined\else
  \let\oldsubparagraph\subparagraph
  \renewcommand{\subparagraph}{
    \@ifstar
      \xxxSubParagraphStar
      \xxxSubParagraphNoStar
  }
  \newcommand{\xxxSubParagraphStar}[1]{\oldsubparagraph*{#1}\mbox{}}
  \newcommand{\xxxSubParagraphNoStar}[1]{\oldsubparagraph{#1}\mbox{}}
\fi
\makeatother


\usepackage{longtable,booktabs,array}
\newcounter{none} % for unnumbered tables
\usepackage{calc} % for calculating minipage widths
% Correct order of tables after \paragraph or \subparagraph
\usepackage{etoolbox}
\makeatletter
\patchcmd\longtable{\par}{\if@noskipsec\mbox{}\fi\par}{}{}
\makeatother
% Allow footnotes in longtable head/foot
\IfFileExists{footnotehyper.sty}{\usepackage{footnotehyper}}{\usepackage{footnote}}
\makesavenoteenv{longtable}
\usepackage{graphicx}
\makeatletter
\newsavebox\pandoc@box
\newcommand*\pandocbounded[1]{% scales image to fit in text height/width
  \sbox\pandoc@box{#1}%
  \Gscale@div\@tempa{\textheight}{\dimexpr\ht\pandoc@box+\dp\pandoc@box\relax}%
  \Gscale@div\@tempb{\linewidth}{\wd\pandoc@box}%
  \ifdim\@tempb\p@<\@tempa\p@\let\@tempa\@tempb\fi% select the smaller of both
  \ifdim\@tempa\p@<\p@\scalebox{\@tempa}{\usebox\pandoc@box}%
  \else\usebox{\pandoc@box}%
  \fi%
}
% Set default figure placement to htbp
\def\fps@figure{htbp}
\makeatother



\ifLuaTeX
\usepackage[bidi=basic,shorthands=off]{babel}
\else
\usepackage[bidi=default,shorthands=off]{babel}
\fi
\ifLuaTeX
  \usepackage{selnolig} % disable illegal ligatures
\fi


\setlength{\emergencystretch}{3em} % prevent overfull lines

\providecommand{\tightlist}{%
  \setlength{\itemsep}{0pt}\setlength{\parskip}{0pt}}



 


\usepackage{enumitem}
\usepackage{booktabs}
\usepackage{tikz}
\setlength{\parskip}{6pt}
\setlength{\parindent}{0pt}
\makeatletter
\@ifpackageloaded{caption}{}{\usepackage{caption}}
\AtBeginDocument{%
\ifdefined\contentsname
  \renewcommand*\contentsname{Table of contents}
\else
  \newcommand\contentsname{Table of contents}
\fi
\ifdefined\listfigurename
  \renewcommand*\listfigurename{List of Figures}
\else
  \newcommand\listfigurename{List of Figures}
\fi
\ifdefined\listtablename
  \renewcommand*\listtablename{List of Tables}
\else
  \newcommand\listtablename{List of Tables}
\fi
\ifdefined\figurename
  \renewcommand*\figurename{Figure}
\else
  \newcommand\figurename{Figure}
\fi
\ifdefined\tablename
  \renewcommand*\tablename{Table}
\else
  \newcommand\tablename{Table}
\fi
}
\@ifpackageloaded{float}{}{\usepackage{float}}
\floatstyle{ruled}
\@ifundefined{c@chapter}{\newfloat{codelisting}{h}{lop}}{\newfloat{codelisting}{h}{lop}[chapter]}
\floatname{codelisting}{Listing}
\newcommand*\listoflistings{\listof{codelisting}{List of Listings}}
\makeatother
\makeatletter
\makeatother
\makeatletter
\@ifpackageloaded{caption}{}{\usepackage{caption}}
\@ifpackageloaded{subcaption}{}{\usepackage{subcaption}}
\makeatother
\usepackage{bookmark}
\IfFileExists{xurl.sty}{\usepackage{xurl}}{} % add URL line breaks if available
\urlstyle{same}
\hypersetup{
  pdftitle={Practice Exam},
  pdfauthor={Philosophy 305: Introduction to Formal Methods},
  pdflang={en},
  colorlinks=true,
  linkcolor={blue},
  filecolor={Maroon},
  citecolor={Blue},
  urlcolor={Blue},
  pdfcreator={LaTeX via pandoc}}


\title{Practice Exam}
\author{Philosophy 305: Introduction to Formal Methods}
\date{2026-04-01}
\begin{document}
\maketitle


\textbf{Time}: This practice exam is designed to be similar too, or
maybe somewhat harder than, the final.

\textbf{Calculators}: permitted.

\textbf{Notes}: Will not be permitted on the final, but obviously
consult your notes in reviewing.

\bigskip
\hrule
\bigskip

\newpage

\section{Question 1: Probability and Conditional
Probability}\label{question-1-probability-and-conditional-probability}

A payment processing company uses an automated system to flag
potentially fraudulent transactions. From historical data, 2\% of all
transactions are fraudulent. The system correctly flags 90\% of
fraudulent transactions. However, it also incorrectly flags 5\% of
legitimate transactions.

\textbf{(a)}: What is the probability that a randomly selected
transaction gets flagged by the system?

\vspace{5cm}

\textbf{(b)}: A transaction has been flagged. What is the probability
that it is actually fraudulent?

\vspace{5cm}

\textbf{(c)}: Is the event ``the transaction is flagged'' independent of
the event ``the transaction is fraudulent''? Justify your answer using
the definition of independence.

\newpage

\section{Question 2: Probability Trees and
Independence}\label{question-2-probability-trees-and-independence}

A student takes a two-part qualifying exam. The probability of passing
Part 1 is 0.7. If they pass Part 1, the probability of passing Part 2 is
0.8 (a confidence boost). If they fail Part 1, the probability of
passing Part 2 drops to 0.3.

\textbf{(a)}: What is the probability that the student passes both
parts?

\vspace{5cm}

\textbf{(b)}: What is the probability that the student passes exactly
one of the two parts?

\newpage

\section{Question 3: Expected
Utility}\label{question-3-expected-utility}

A small business owner must choose between two expansion plans before
learning the economic conditions for the coming year. There are three
possible scenarios:

{\def\LTcaptype{none} % do not increment counter
\begin{longtable}[]{@{}
  >{\raggedleft\arraybackslash}p{(\linewidth - 6\tabcolsep) * \real{0.1449}}
  >{\centering\arraybackslash}p{(\linewidth - 6\tabcolsep) * \real{0.3043}}
  >{\centering\arraybackslash}p{(\linewidth - 6\tabcolsep) * \real{0.2754}}
  >{\centering\arraybackslash}p{(\linewidth - 6\tabcolsep) * \real{0.2754}}@{}}
\toprule\noalign{}
\begin{minipage}[b]{\linewidth}\raggedleft
\end{minipage} & \begin{minipage}[b]{\linewidth}\centering
Downturn (Pr = 0.3)
\end{minipage} & \begin{minipage}[b]{\linewidth}\centering
Stable (Pr = 0.4)
\end{minipage} & \begin{minipage}[b]{\linewidth}\centering
Growth (Pr = 0.3)
\end{minipage} \\
\midrule\noalign{}
\endhead
\bottomrule\noalign{}
\endlastfoot
Plan X & -90 & 80 & 100 \\
Plan Y & 10 & 30 & 50 \\
\end{longtable}
}

Profits are in thousands of dollars.

\textbf{(a)}: Which plan has higher expected monetary value? Show the
calculation for each.

\vspace{4cm}

\textbf{(b)}: Does either plan dominate the other (strictly or weakly)?
Explain briefly.

\vspace{4cm}

\textbf{(c)}: The owner currently has \$100,000 in savings and has a
utility function
\(U(\text{total wealth}) = \sqrt{\text{total wealth in thousands}}\).
So, for example, if the owner chooses Plan X and a Downturn occurs,
total wealth is \$10,000 and utility is \(\sqrt{10}\). Which plan should
the owner choose according to expected utility? Does this differ from
your answer in part (a)?

\newpage

\section{Question 4: Strategic-Form
Games}\label{question-4-strategic-form-games}

Consider the following two-player game, where payoffs are listed as
(Row, Column):

{\def\LTcaptype{none} % do not increment counter
\begin{longtable}[]{@{}rcccc@{}}
\toprule\noalign{}
& A & B & C & D \\
\midrule\noalign{}
\endhead
\bottomrule\noalign{}
\endlastfoot
\textbf{T} & (4, 2) & (2, 1) & (3, 3) & (1, 0) \\
\textbf{M} & (2, 3) & (4, 0) & (1, 2) & (3, 1) \\
\textbf{B} & (1, 1) & (3, 2) & (2, 4) & (0, 0) \\
\end{longtable}
}

\textbf{(a)}: In the original game, does any strategy strictly dominate
another for either player? State all dominance relations you find.

\vspace{4cm}

\textbf{(b)}: Use iterated elimination of strictly dominated strategies
to simplify the game. Show each round of elimination clearly, stating
which strategy is eliminated and why.

\vspace{6cm}

\textbf{(c)}: Find all pure-strategy Nash equilibria of the original
game. Use the best-response method and show your work.

\newpage

\section{Question 5: Dynamic Games and Backward
Induction}\label{question-5-dynamic-games-and-backward-induction}

Consider the following game in extensive form. Player 1 moves first,
choosing \textbf{A} or \textbf{B}. Player 2 observes Player 1's choice
and responds. If Player 1 chose A, Player 2 chooses \textbf{C} or
\textbf{D}. If Player 1 chose B, Player 2 chooses \textbf{E} or
\textbf{F}. If Player 2 chooses F, then Player 1 gets a second move,
choosing \textbf{G} or \textbf{H}. Payoffs are (Player 1, Player 2):

\begin{center}
\begin{tikzpicture}[
  every node/.style={font=\small},
  dot/.style={circle, fill=black, inner sep=1.5pt},
  payoff/.style={font=\small}
]
% Player 1 root (open circle)
\node[circle, draw, inner sep=1.5pt, label=left:{\textbf{P1}}] (root) at (0,0) {};

% A branch
\node[dot, label=above left:{\textbf{P2}}] (p2a) at (3, 2) {};
\draw (root) -- node[above left] {$A$} (p2a);

\node[payoff] (ac) at (6, 3) {$(5,\, 3)$};
\node[payoff] (ad) at (6, 1) {$(2,\, 1)$};
\draw (p2a) -- node[above] {$C$} (ac);
\draw (p2a) -- node[below] {$D$} (ad);

% B branch
\node[dot, label=below left:{\textbf{P2}}] (p2b) at (3, -2) {};
\draw (root) -- node[below left] {$B$} (p2b);

\node[payoff] (be) at (6, -1) {$(3,\, 4)$};
\draw (p2b) -- node[above] {$E$} (be);

% F sub-branch
\node[dot, label=below left:{\textbf{P1}}] (p1b) at (6, -3) {};
\draw (p2b) -- node[below] {$F$} (p1b);

\node[payoff] (bg) at (9, -2) {$(6,\, 0)$};
\node[payoff] (bh) at (9, -4) {$(1,\, 5)$};
\draw (p1b) -- node[above] {$G$} (bg);
\draw (p1b) -- node[below] {$H$} (bh);
\end{tikzpicture}
\end{center}

\textbf{(a)}: How many strategies does each player have? List them all.
(Remember: a strategy specifies a choice at every decision node where
that player moves, including nodes that may not be reached.)

\vspace{5cm}

\textbf{(b)}: Solve this game using backward induction. State clearly
what happens at each step and what the final outcome is.

\vspace{5cm}

\textbf{(c)}: Player 2 would prefer the outcome (3, 4) to the
backward-induction outcome. Suppose Player 2 threatens to play D (rather
than C) after A, hoping to push Player 1 toward B. Is this threat
credible? Explain.

\vspace{4cm}

\textbf{(d)}: Write out the strategic (normal) form of this game. That
is, construct the payoff table with Player 1's strategies as rows and
Player 2's strategies as columns.

\newpage

\section{Question 6: Mixed
Strategies}\label{question-6-mixed-strategies}

Consider the following game:

{\def\LTcaptype{none} % do not increment counter
\begin{longtable}[]{@{}rcc@{}}
\toprule\noalign{}
& Left & Right \\
\midrule\noalign{}
\endhead
\bottomrule\noalign{}
\endlastfoot
\textbf{Up} & (3, 1) & (1, 2) \\
\textbf{Down} & (1, 3) & (2, 0) \\
\end{longtable}
}

This game has no pure-strategy Nash equilibrium. Find the mixed-strategy
Nash equilibrium. That is, find the probability with which Row plays Up
and the probability with which Column plays Left such that neither
player can improve their expected payoff. Also calculate each player's
expected payoff in the equilibrium.

\vspace{10cm}

\bigskip
\hrule
\bigskip
\begin{center}
\textbf{END OF PRACTICE EXAM}
\end{center}




\end{document}
